\section{Related work and the benchmark landscape}
\label{sec:related}

\subsection{Dynamic Bayesian networks and structure learning}
\label{sec:related_dbn}

Dynamic Bayesian networks generalize hidden Markov models and linear state-space
models by writing a temporal process as a factored graphical model over a
repeated time slice \citep{dean1989model, murphy2002dbn}. The representation is
attractive precisely when the temporal dependence has to remain readable, since
each edge names a variable, a lag and a target. Structure learning for the
dynamic case inherits the machinery of the static case. Score-based search
maximizes a decomposable criterion, most often the Bayesian information criterion
or a Bayesian Dirichlet score \citep{heckerman1995learning}; constraint-based
search recovers a skeleton from conditional independence tests
\citep{spirtes2000causation}; greedy equivalence search operates on equivalence
classes \citep{chickering2002optimal}; and hybrid methods combine a local
neighborhood step with a scored search \citep{tsamardinos2006max}. The tree
special case has an exact maximizer \citep{chow1968approximating}, the
time-series literature contributes lag selection read as model selection
\citep{granger1969investigating}, and continuous-optimization formulations
replace the combinatorial search with a smooth acyclicity penalty and an
autoregressive design \citep{pamfil2020dynotears}. Open implementations of most
of these are available in a single library \citep{ankan2015pgmpy, ankan2024pgmpy}.

Two features of that literature matter here. Every one of these criteria is
computed from the observed variables alone, so none of them conditions on a
latent state while scoring a candidate parent. And the flagship survey of
Bayesian network structure learning covers 74 algorithms and states explicitly
that structure learning from time-series data is outside its scope, deferring to
a separate review \citep{kitson2023bnsurvey, nogueira2022cdtools}. The dynamic
case is a distinct literature with distinct tooling, and the present work sits
inside it.

\subsection{Non-stationary and time-varying variants}
\label{sec:related_nonstationary}

Classical dynamic Bayesian network learning assumes a stationary generating
process, which long behavioral horizons violate. The response splits into two
premises. In the first, the graph itself changes, and learning means recovering a
sequence of structures together with the times at which they change
\citep{robinson2010learning, lebre2010artiva}. In the second, the graph is shared
and only the parameters change across segments, which is the premise of the
allocation and coupling models of Grzegorczyk and Husmeier and of their
successors \citep{grzegorczyk2011nhdbncontinuous, dondelinger2013edison}. A
parallel line treats the parameters as varying smoothly rather than in segments,
estimating a time-varying network by kernel-weighted or penalized regression
\citep{song2009tvdbn, kolar2010estimating}. More recent work makes the regime an
explicit latent variable to be inferred jointly with per-regime structure
\citep{saggioro2020rpcmci, rahmani2025castor}.

The n\(^2\)DBN belongs to the second premise. The lagged graph at a scale is
shared across the horizon and the regime \(S_t^{(k)}\) modulates the emission and
allocation parameters, with the regime chain at a finer scale conditioned on its
parent-scale regime so that drift is coherent across resolutions. What separates
this model from the cited work is that the regime is one of two latent
mechanisms rather than the only one, and that the scales are coupled by an exact
aggregation identity rather than by a shared prior.

There is a specific hazard in this setting that the non-stationary literature has
already named. A persistent latent state raises the level of every variable it
drives for as long as it lasts, so each variable becomes positively
autocorrelated even with an empty parent set, and a criterion that conditions on
nothing attributes that autocorrelation to a lagged self edge
\citep{runge2018pitfalls}. This is not a defect of any particular algorithm. It
is a property of the hypothesis class those algorithms search, and it is the
mechanism the comparison in this paper measures.

\subsection{Temporal abstraction and hierarchical time series}
\label{sec:related_abstraction}

Reasoning at several temporal granularities rather than only over instantaneous
slices is an old concern in the Bayesian network literature, where temporal
abstraction supplies the vocabulary of intervals and episodes that a slice-based
model lacks \citep{burns2003temporal}. The forecasting literature supplies the
other half of the picture. A hierarchical time series carries an exact and known
summing map from the bottom level to the aggregates, and independent forecasts of
the levels are almost never coherent with it, so a reconciliation step projects
them back onto the space that satisfies the summing constraints
\citep{athanasopoulos2009tourism, wickramasuriya2019mint}. The summing map of
that literature is the direct analogue of \cref{eq:intro_reconciliation}.

The n\(^2\)DBN takes the constraint from the second literature and the episode
semantics from the first, and differs from both in where the constraint is
imposed. Reconciliation in hierarchical forecasting is a projection applied after
independent models have been fitted, and it is optimal under an assumption about
the base forecast errors. Here the constraint is a property of the generative
model, since a finer scale exists only as an allocation of its parent total, so
there is no incoherent intermediate object to project and no reconciliation
estimator to specify. The cost of that choice is that the fine scale is not free.
It competes for a fixed total, which turns out to have consequences the
limitations record.

\subsection{The benchmark landscape}
\label{sec:related_benchmarks}

A model whose central property is exact cross-scale aggregation invites an
obvious question. Why not score it on a public benchmark rather than on a
generator written for the purpose? We surveyed the datasets available for scoring
a dynamic Bayesian network learner and found a specific and citable gap.

Restricting attention to resources that carry both a released ground-truth graph
and a genuine time series, which is the set on which a structure learner can be
scored at all, leaves twelve. They are the DREAM4 and DREAM3 in-silico network
challenges \citep{marbach2010revealing, prill2010dream3} and the GeneNetWeaver
simulator behind them \citep{schaffter2011gnw}; the TimeGraph synthetic suite
\citep{ferdous2025timegraph}; CausalDynamics \citep{herdeanu2025causaldynamics};
CausalTime \citep{cheng2024causaltime}; the Causal Chambers
\citep{gamella2025causalchambers}; CausalRivers \citep{stein2025causalrivers};
the IRMA yeast circuit \citep{cantone2009irma}; the NetSim simulated fMRI
networks \citep{smith2011netsim}; FinanceCPT \citep{nauta2019tcdf}; and
CIPCaD-Bench \citep{menegozzo2022cipcad}. Two kinds of resource frequently
listed alongside them do not belong, since SynTReN generates steady-state data
rather than a time series, and the pseudotime benchmarks order cells along a
differentiation trajectory over a network that does not change, so what varies
across their samples is not time.

Not one of the twelve combines the three properties this model is about, namely
event counts, exact cross-scale aggregation and a latent regime. The gap is not
that any one of the three is unattainable. A regime change with something to
score against is present in five of the twelve, and count series that aggregate
exactly over a known hierarchy are routine in forecasting repositories. What no
resource does is release a ground-truth graph alongside all three at once.

DREAM4 is the most used quantitative benchmark in the field, with five networks
of ten genes and five of a hundred, replicated time series of twenty-one points
each, and a released directed gold standard; DREAM3 is the same construction at
ten, fifty and a hundred genes. The observations are continuous expression levels
rather than counts, the design is single resolution, and the perturbation applied
and then removed is a stimulus rather than a structural change, so the process is
stationary. GeneNetWeaver is the simulator that produced them and generates the
same kind of data at chosen sizes, with the same three absences.

TimeGraph is the closest synthetic match on the non-stationarity axis, since
non-stationarity is an explicit generator axis alongside trends, seasonality,
irregular sampling and missingness, and it ships the structural equations rather
than only the graph. Its series are four to eight continuous variables over 500
to 3000 steps at a single resolution, so it has a changing generator but no
counts and no aggregation.

The Causal Chambers are the strongest real-data item, and the only one giving a
real physical system, a known graph, known intervention targets and marked
changepoints at once. The regime-jump datasets run 320{,}000 observations in
thirty-two regimes of ten thousand, with an intervention drawn from an
eleven-variable pool at seventy percent of the regime boundaries. The observed
quantities are sensor readings rather than event counts, the regime is an applied
intervention rather than a latent state to be inferred, and there is no nesting,
so any hierarchy would be imposed by the analyst rather than intrinsic.

CausalRivers is the largest real-world time-series causal benchmark, with
hundreds of stations at fifteen-minute resolution over several years and a
released station graph, and its flood product contains a genuine distributional
shift. The shift is real, it is not labeled as a regime path, the measurements
are water levels rather than counts, and nothing aggregates exactly.

IRMA is the only classical real dataset with both a known network and an
experimentally imposed regime switch, since the growth medium is changed inside
the series. It is five genes, eight links and series of sixteen or twenty points,
which supports a demonstration and not a power analysis, and its readings are
again continuous and single-scale.

CausalDynamics is the largest synthetic resource, with more than fourteen
thousand graphs and more than fifty million samples drawn from coupled
differential equations and two idealized climate models, and its adjacency is
derived analytically from the governing equations. The trajectories are chaotic
or stochastic rather than regime-switching by design, and they are continuous and
single-scale.

The remaining four complete the picture without changing it. NetSim and
FinanceCPT are stationary simulations on fixed graphs, the second usefully
labeling its ground truth with the lag. CausalTime fits a deep model to real data
and emits synthetic series that keep real dynamics and carry a graph by
construction, which is the right instrument against the charge that fixed
synthetic benchmarks are too easy, but the graph is generated rather than
externally validated and none of the three properties is present. CIPCaD-Bench
supplies an industrial process with faults injected mid-run and a graph derived
from a plant flowsheet, which is the closest thing to a regime among the
industrial data, again without counts or nesting.

The tempting alternative is a static repository. The bnlearn Bayesian Network
Repository distributes fully parameterized networks with an exact directed graph
and exact conditional probability tables \citep{scutari2010bnlearn}, and a
package vendors them so that an experiment runs without network access
\citep{ankan2024pgmpy}. It has no temporal section of any kind. Turning one of
its networks into a two-slice temporal network means mapping each edge onto a
lag-1 edge and sampling, which is what the field already does when it needs a
dynamic benchmark, with published methodology and published precedent
\citep{trabelsi2013benchmarking, ouyang2023divideconquer}. The construction is
sound and this work uses it, but it manufactures the time series it needs. The
tiled network has no nesting, no counts that aggregate and no latent regime, so
three of the four requirements of \cref{sec:intro_requirements} are absent by
construction. There is also no dynamic counterpart to the repository itself. The
2025 expansion of the static collection adds a further large set of networks
drawn from the literature and no dynamic ones \citep{leonelli2025bnrep}, and the
reference reproducible
workflow for benchmarking structure learning has no temporal graph generator, no
time-series data module and no dynamic algorithm entries
\citep{rios2025benchpress}.

The conclusion is narrow and it is the justification for this work shipping its
own benchmark. Multiscale reconciliation cannot be scored on any released
benchmark, because no released benchmark contains it. What can be scored
externally is the structure learner, the local score and the table estimator on
flat discrete data, which is what the tiled networks are used for here, and the
resulting comparison is reported alongside the generated one rather than in place
of it.

\subsection{Released parameters, not only released graphs}
\label{sec:related_parameters}

A recovery claim about parameters needs more than a released graph, since the
generating conditional distributions have to be known exactly for the comparison
to mean anything. Most of the twelve release a graph and nothing else. Four
resources release true parameters. GeneNetWeaver attaches kinetic models of
transcription and translation to the subnetworks it extracts, so a generated
network carries its own parameters by construction. TimeGraph releases the
structural equations and the full generating metadata alongside each graph.
CausalDynamics derives its adjacency from governing equations that are themselves
released. And the tiled bnlearn networks carry exact conditional probability
tables, which is why the external check in this paper can score an estimated
table against a published one rather than only counting edges. The real systems
in the list, namely the Causal Chambers, CausalRivers and IRMA, have no true
parameters to release, and the DREAM challenges release edge lists rather than
the kinetics of the specific challenge networks.

Two further hazards of benchmark design deserve mention, since both would be
raised against any evaluation of this kind. A synthetic generator for a two-slice
network can induce an ordering of the variables by marginal variance that
continuous-optimization learners exploit for the wrong reason, so standardization
and a report of the residual sortability are both needed
\citep{reisach2021varsortability}. And structural Hamming distance is defined
over partially directed graphs in its original form \citep{tsamardinos2006max}
while several widely used libraries compute it on the directed adjacency, so a
paper has to state which convention it used before its distances can be compared
with anyone else's.

Finally, a falsification test that requires no ground truth exists and transfers
to real logs. The self-compatibility check compares graphs learned from different
variable subsets and asks whether they agree where they overlap
\citep{faller2024selfcompatibility}. It is necessary rather than sufficient, and
it is the only structural check available where the true graph is unknown.
